% =====================================================================
% Phase 1 — LangChain LCEL : 5 concepts essentiels
% Document unifié regroupant ChatPromptTemplate, sortie structurée Pydantic,
% composition LCEL + .partial(), RunnableParallel, et streaming.
% Compilation : pdflatex -shell-escape phase1.tex (deux fois pour la TOC)
% Pré-requis  : pygmentize (pip install Pygments) + paquet TeX "minted"
% =====================================================================
\documentclass[11pt,a4paper]{article}

% ---------- Encodage & langue ----------
\usepackage[utf8]{inputenc}
\usepackage[T1]{fontenc}
\usepackage[french]{babel}
\usepackage{lmodern}
\usepackage{microtype}

% ---------- Mise en page ----------
\usepackage[a4paper,margin=2.2cm]{geometry}
\usepackage{parskip}
\usepackage{xcolor}
\usepackage{enumitem}
\usepackage{booktabs}
\usepackage{hyperref}
\hypersetup{
  colorlinks=true,
  linkcolor=blue!60!black,
  urlcolor=blue!60!black,
  citecolor=blue!60!black,
  bookmarksopen=true,
  bookmarksopenlevel=2,
}

% ---------- Boîtes & callouts ----------
\usepackage[most]{tcolorbox}
\tcbset{
  enhanced,
  colback=blue!3,
  colframe=blue!50!black,
  fonttitle=\bfseries,
  arc=2pt, boxrule=0.6pt,
}

% ---------- Code avec minted ----------
\usepackage{minted}
\setminted{
  fontsize=\small,
  linenos=true,
  numbersep=6pt,
  frame=lines,
  framesep=2mm,
  bgcolor=gray!4,
  breaklines=true,
  breakanywhere=true,
  python3=true,
  tabsize=4,
}
\newminted{python}{}

% ---------- TikZ (union des librairies des 5 concepts) ----------
\usepackage{tikz}
\usetikzlibrary{
  arrows.meta, positioning, shapes.geometric, shapes.misc,
  fit, backgrounds, calc, trees,
  decorations.pathreplacing, decorations.markings,
}

\tikzset{
  node distance=8mm and 12mm,
  block/.style={
    rectangle, rounded corners=2pt, draw=blue!60!black, fill=blue!8,
    minimum width=42mm, minimum height=9mm, align=center, font=\small
  },
  data/.style={
    rectangle, rounded corners=1pt, draw=gray!60, fill=gray!10,
    minimum width=42mm, minimum height=8mm, align=center, font=\footnotesize\ttfamily
  },
  llm/.style={
    rectangle, rounded corners=2pt, draw=red!60!black, fill=red!8,
    minimum width=42mm, minimum height=9mm, align=center, font=\small
  },
  msg/.style={
    rectangle, rounded corners=1pt, draw=purple!60!black, fill=purple!8,
    minimum width=56mm, minimum height=8mm, align=left, font=\scriptsize\ttfamily
  },
  schema/.style={
    rectangle, rounded corners=2pt, draw=teal!60!black, fill=teal!8,
    minimum width=48mm, minimum height=9mm, align=center, font=\small
  },
  pipeop/.style={
    circle, draw=blue!60!black, fill=blue!15, minimum size=6mm, font=\bfseries\small
  },
  parallel/.style={
    rectangle, rounded corners=2pt, draw=violet!60!black, fill=violet!8,
    minimum width=48mm, minimum height=9mm, align=center, font=\small
  },
  token/.style={
    rectangle, rounded corners=1pt, draw=teal!60!black, fill=teal!10,
    minimum height=6mm, font=\scriptsize\ttfamily, inner sep=2pt
  },
  flow/.style={-{Stealth[length=2.2mm]}, thick, draw=blue!60!black},
  flowdata/.style={-{Stealth[length=2.2mm]}, thick, draw=gray!60, dashed},
  label/.style={font=\scriptsize\itshape, fill=white, inner sep=1pt},
  timearrow/.style={-{Stealth[length=2.5mm]}, very thick, draw=black!70},
  bar/.style={rectangle, draw=blue!60!black, fill=blue!20, minimum height=4mm, font=\scriptsize},
}

% ---------- Titre ----------
\title{\textbf{Phase 1 — LangChain LCEL}\\
\large Le moteur d'explication structurée : cinq concepts essentiels}
\author{Cours PyTutor2}
\date{}

\begin{document}
\maketitle

\begin{abstract}
\noindent
Ce document couvre les cinq briques fondamentales de la phase 1 du cours
PyTutor2~: les \texttt{ChatPromptTemplate}, la sortie typée avec Pydantic
et \texttt{with\_structured\_output}, la composition LCEL avec \texttt{|}
et \texttt{.partial()}, l'exécution parallèle avec
\texttt{RunnableParallel}, et le streaming token par token avec
\texttt{.stream()}. Chaque concept est traité avec un énoncé, une
architecture illustrée en TikZ, le code complet, des explications pas à pas
et la vérification dans LangSmith.
\end{abstract}

\tableofcontents
\bigskip

% #####################################################################
% #####################################################################
\section{\texttt{ChatPromptTemplate} : structurer un prompt en messages paramétrés}
% #####################################################################
% #####################################################################

\begin{tcolorbox}[title=Exercice 1. \texttt{ChatPromptTemplate}]
On veut produire un assistant pédagogique Python dont le \textbf{ton et la
profondeur} s'adaptent au niveau de l'élève (débutant, intermédiaire, avancé).
Le concept à expliquer change à chaque requête, le niveau aussi.

\textbf{Contrainte}~: ne pas écrire trois prompts différents, ni concaténer des
strings à la main. On veut \emph{un seul} template réutilisable, avec deux
variables \texttt{\{level\}} et \texttt{\{concept\}}, qui sera ensuite consommé
par n'importe quel modèle compatible LangChain.

\bigskip
\noindent\textbf{Objectifs pédagogiques}
\begin{itemize}[leftmargin=*]
  \item Comprendre pourquoi un prompt moderne est une \textbf{liste de messages typés},
        pas une chaîne de caractères.
  \item Maîtriser la séparation \texttt{system} / \texttt{human} et son rôle pour
        la qualité des réponses.
  \item Utiliser les \textbf{variables} \texttt{\{nom\}} d'un template et savoir
        ce qu'elles deviennent au moment du rendu.
\end{itemize}
\end{tcolorbox}

\subsection{Solution}

\subsubsection{Architecture}

\begin{center}
\begin{tikzpicture}
  \node[data] (in) {\{"level": "débutant",\\\ "concept": "lambda"\}};
  \node[block, below=of in] (tpl) {ChatPromptTemplate\\\texttt{from\_messages([...])}};
  \node[msg, below=of tpl, text width=72mm] (msgs)
    {SystemMessage("Tu es un assistant\dots\\\hspace*{1em}Niveau actuel : débutant\dots")\\
     HumanMessage("Explique : lambda")};
  \node[block, below=of msgs] (model) {Modèle de chat\\(ChatAnthropic, etc.)};

  \draw[flowdata] (in)   -- (tpl);
  \draw[flow]     (tpl)  -- (msgs) node[label, midway, right]{\texttt{.invoke(dict)}};
  \draw[flowdata] (msgs) -- (model);
\end{tikzpicture}
\end{center}

Un \texttt{ChatPromptTemplate} est essentiellement une \textbf{fonction}
\texttt{dict $\to$ List[BaseMessage]}. Les variables \texttt{\{level\}} et
\texttt{\{concept\}} sont substituées dans le contenu de chaque message au
moment du rendu.

\subsubsection{Le code complet}

\begin{pythoncode}
from langchain_core.prompts import ChatPromptTemplate

EXPLAIN_PROMPT = ChatPromptTemplate.from_messages([
    (
        "system",
        "Tu es un assistant pédagogique pour des élèves apprenant Python. "
        "Tu adaptes systématiquement ton vocabulaire et tes exemples au "
        "niveau de l'élève. Niveau actuel : {level}.\n\n"
        "Règles selon le niveau :\n"
        "  - 'débutant' : pas de jargon, exemples ultra-simples, "
        "analogies du quotidien, pas plus de 10 lignes de code.\n"
        "  - 'intermédiaire' : tu peux supposer la syntaxe de base "
        "connue, montre les usages idiomatiques (PEP 8, type hints).\n"
        "  - 'avancé' : montre les détails du modèle objet de Python, "
        "les subtilités, les cas limites, les implications de performance."
    ),
    ("human", "Explique le concept suivant : {concept}"),
])
\end{pythoncode}

\subsubsection{Explications pas à pas}

\subsubsection*{Étape 1 — Pourquoi une \emph{liste de messages} et pas un string}

Les modèles de chat modernes (Claude, GPT-4, Gemini, etc.) ne reçoivent pas
un prompt monolithique~: ils reçoivent une \textbf{conversation structurée}
sous forme de messages typés. Chaque message a un \texttt{role}
(\texttt{system}, \texttt{user}/\texttt{human}, \texttt{assistant}) et un
contenu. L'API du fournisseur les traite différemment~:

\begin{itemize}[leftmargin=*]
  \item Le \texttt{system} fixe la \textbf{personnalité} et les règles du
        jeu, et reçoit le plus de poids dans le respect des instructions.
  \item Les \texttt{human} sont les requêtes de l'utilisateur.
  \item Les \texttt{assistant} sont les réponses précédentes du modèle
        (utile pour du multi-tour, on les verra plus tard).
\end{itemize}

\begin{quote}\itshape
Mettre les \emph{règles} dans le \texttt{system} et la \emph{requête} dans
le \texttt{human} n'est pas cosmétique~: les LLM modernes obéissent
nettement mieux à une instruction présente dans le message système qu'à la
même instruction noyée dans le message utilisateur. Si le ton de ton
assistant dérive, c'est souvent que tu as mis tes règles au mauvais endroit.
\end{quote}

\subsubsection*{Étape 2 — La syntaxe \texttt{from\_messages([...])}}

\texttt{from\_messages} accepte une liste de \textbf{tuples} de la forme
\texttt{(role, content)}~:

\begin{center}
\begin{tabular}{ll}
\texttt{("system", "\dots")}     & devient un \texttt{SystemMessagePromptTemplate} \\
\texttt{("human",  "\dots")}     & devient un \texttt{HumanMessagePromptTemplate} \\
\texttt{("ai",     "\dots")}     & devient un \texttt{AIMessagePromptTemplate} \\
\texttt{("placeholder", "\{x\}")} & insère une liste de messages stockée dans \texttt{x} \\
\end{tabular}
\end{center}

LangChain transforme automatiquement chaque tuple en un objet
\texttt{PromptTemplate} de la bonne classe, et l'ensemble compose un
\texttt{ChatPromptTemplate}.

\subsubsection*{Étape 3 — Les variables \texttt{\{nom\}}}

Toutes les accolades non-doublées dans une chaîne de message sont
considérées comme des \textbf{variables}~:

\begin{itemize}[leftmargin=*]
  \item \texttt{\{concept\}} dans le \texttt{human} attend une clé
        \texttt{concept} dans le dict d'entrée.
  \item \texttt{\{level\}} dans le \texttt{system} attend une clé
        \texttt{level}.
  \item Si une clé manque au moment du rendu, LangChain lève une
        \texttt{KeyError} explicite.
\end{itemize}

\begin{quote}\itshape
Si tu veux écrire une \emph{vraie} accolade dans le prompt (par exemple,
demander au modèle de répondre en JSON), il faut la \textbf{doubler}~:
\texttt{\{\{} pour un \texttt{\{} littéral, \texttt{\}\}} pour un \texttt{\}}.
C'est le piège classique qui transforme un message en
\texttt{KeyError: 'nom\_de\_clé\_inattendu'} dès qu'on parle de structures
JSON dans le prompt.
\end{quote}

\subsubsection*{Étape 4 — Le rendu via \texttt{.invoke()} ou \texttt{.format\_messages()}}

\begin{center}
\resizebox{\linewidth}{!}{%
\begin{tikzpicture}[node distance=6mm and 14mm]
  \node[data] (in) {\{"level": "débutant",\\\ "concept": "lambda"\}};
  \node[block, right=of in] (tpl) {EXPLAIN\_PROMPT};
  \node[msg, right=of tpl, text width=66mm] (out)
    {[SystemMessage("\dots Niveau : débutant\dots"),\\\ HumanMessage("Explique : lambda")]};

  \draw[flowdata] (in)  -- (tpl);
  \draw[flow]     (tpl) -- (out) node[label, midway, above]{\texttt{.invoke(\dots)}};
\end{tikzpicture}%
}
\end{center}

Deux façons d'appeler le template~:

\begin{enumerate}[leftmargin=*]
  \item \texttt{EXPLAIN\_PROMPT.invoke(\{"level": "débutant", "concept": "lambda"\})}
        renvoie un \texttt{ChatPromptValue} (qu'on passe ensuite au modèle).
        C'est la voie LCEL, utilisée quand le template est dans une chaîne
        \texttt{prompt | model | \dots}
  \item \texttt{EXPLAIN\_PROMPT.format\_messages(\dots)} renvoie directement
        la liste de \texttt{BaseMessage}. Pratique pour \textbf{déboguer} ou
        pour passer les messages à une API tierce.
\end{enumerate}

\begin{quote}\itshape
\texttt{.format\_messages(\dots)} est le \emph{print} du prompt engineer.
Avant de te demander pourquoi le modèle répond bizarrement, imprime ce que
\textbf{tu lui as réellement envoyé} — souvent l'erreur est là.
\end{quote}

\subsection{Vérification dans LangSmith}

\subsubsection*{Comment voir le prompt rendu}

Quand le template est utilisé dans une chaîne et que LangSmith est activé,
chaque appel au modèle produit un run \texttt{ChatAnthropic} (ou équivalent)
dans lequel l'onglet \emph{Input} contient \textbf{exactement les messages
rendus}~:

\begin{center}
\begin{tikzpicture}[every node/.style={font=\footnotesize\ttfamily, anchor=west}]
  \node                              (a) {RunnableSequence};
  \node[below=1mm of a, xshift=4mm]  (b) {ChatPromptTemplate};
  \node[below=1mm of b, xshift=4mm]  (c) {Input  : \{level: "débutant", concept: "lambda"\}};
  \node[below=1mm of c]              (d) {Output : [SystemMessage(...), HumanMessage(...)]};
  \node[below=1mm of d, xshift=-4mm] (e) {ChatAnthropic};
  \node[below=1mm of e, xshift=4mm]  (f) {Input  : la liste de messages ci-dessus};
\end{tikzpicture}
\end{center}

\subsubsection*{Les trois choses à regarder}

\begin{enumerate}[leftmargin=*]
  \item \textbf{L'onglet \emph{Input} du nœud \texttt{ChatPromptTemplate}}~: tu
        vois le dict d'entrée brut. C'est là que tu repères une variable
        manquante ou mal typée.
  \item \textbf{L'onglet \emph{Output} du même nœud}~: la liste complète des
        messages, variables substituées, prêts à partir au LLM. C'est la
        \textbf{vérité terrain} de ce que le modèle reçoit.
  \item \textbf{L'onglet \emph{Input} du nœud \texttt{ChatAnthropic}}~: en
        principe identique à l'output du template ci-dessus. Si tu y vois
        autre chose, c'est qu'un \texttt{Runnable} intermédiaire a modifié
        les messages.
\end{enumerate}

\begin{quote}\itshape
Change la valeur de \texttt{level} entre \texttt{"débutant"} et
\texttt{"avancé"} pour le \emph{même} concept. Dans LangSmith, tu vois
\textbf{la même structure de messages}, seul le contenu du \texttt{system}
diffère. C'est la visualisation directe de la paramétrisation par template.
\end{quote}

\subsection*{Récapitulatif}
{\itshape\sloppy
\textit{Ce qu'il faut retenir :}
\begin{enumerate}[leftmargin=*]
  \item Un prompt de chat est une \textbf{liste de messages typés}, pas un
        string. Le rôle (\texttt{system}/\texttt{human}/\texttt{ai}) change
        radicalement la façon dont le modèle traite le contenu.
  \item \texttt{ChatPromptTemplate.from\_messages([...])} est la voie
        idiomatique pour construire ce prompt. Les tuples \texttt{(role,
        content)} sont la forme la plus lisible.
  \item Les \textbf{variables} \texttt{\{nom\}} dans le contenu deviennent
        des clés attendues du dict d'entrée. Doubler les accolades pour les
        accolades littérales.
  \item \texttt{prompt.invoke(dict)} renvoie un \texttt{ChatPromptValue}
        pour la chaîne LCEL ; \texttt{prompt.format\_messages(dict)} renvoie
        la liste brute, indispensable pour le débogage.
  \item Mettre les règles dans le \texttt{system}, la requête dans le
        \texttt{human}. C'est l'erreur \#1 qui fait dériver le ton d'un
        assistant.
\end{enumerate}
\par}

% #####################################################################
% #####################################################################
\section{Sortie typée : schémas Pydantic et \texttt{with\_structured\_output}}
% #####################################################################
% #####################################################################

\begin{tcolorbox}[title=Exercice 2. Sortie typée avec Pydantic]
Un assistant pédagogique ne doit pas répondre par un pavé de texte
markdown. On veut un objet manipulable~: une \emph{explication structurée}
contenant l'intuition, un exemple de code, un piège fréquent et une liste
de questions de compréhension. Chaque champ a une sémantique précise,
chaque question a sa réponse attendue.

\textbf{Contrainte}~: la réponse doit être \emph{immédiatement utilisable}
en Python, sans étape de parsing fragile. \texttt{result.intuition} doit
fonctionner ; \texttt{result.comprehension\_questions[0].expected\_answer}
aussi.

\bigskip
\noindent\textbf{Objectifs pédagogiques}
\begin{itemize}[leftmargin=*]
  \item Comprendre pourquoi exiger une \textbf{sortie structurée} d'un LLM
        change la nature du code qu'on écrit autour.
  \item Décrire un schéma de réponse avec \texttt{pydantic.BaseModel} et
        \texttt{Field(description=\dots)}.
  \item Maîtriser \texttt{model.with\_structured\_output(Schema)} et savoir
        ce qui se passe \textbf{sous le capot} (tool-calling).
  \item Reconnaître les cas où il faut \emph{imbriquer} des sous-schémas.
\end{itemize}
\end{tcolorbox}

\subsection{Solution}

\subsubsection{Architecture}

\subsubsection*{Vue d'ensemble}

\begin{center}
\begin{tikzpicture}
  \node[data]                 (in)     {\{"concept": "générateurs",\\\ "level": "intermédiaire"\}};
  \node[block, below=of in]   (prompt) {ChatPromptTemplate};
  \node[schema, right=14mm of prompt] (schema) {Schéma Pydantic\\\texttt{Explanation}};
  \node[llm,   below=of prompt] (model) {Modèle de chat\\\texttt{.with\_structured\_output(\dots)}};
  \node[data, below=of model] (out)    {Explanation(\\\hspace*{1em}intuition=\dots,\\\hspace*{1em}code\_example=\dots,\\\hspace*{1em}pitfall=\dots,\\\hspace*{1em}comprehension\_questions=[\dots])};

  \draw[flowdata] (in)     -- (prompt);
  \draw[flow]     (prompt) -- (model);
  \draw[flow]     (schema) -- (model) node[label, midway, above]{contraint};
  \draw[flow]     (model)  -- (out)   node[label, midway, right]{objet Python};
\end{tikzpicture}
\end{center}

Le schéma Pydantic est \textbf{injecté dans l'appel} au modèle. Le SDK
Anthropic l'expose comme un \emph{tool} que Claude doit appeler ; la
réponse en JSON est ensuite désérialisée vers un vrai objet Python.

\subsubsection{Le code complet}

\begin{pythoncode}
from pydantic import BaseModel, Field

class ComprehensionQuestion(BaseModel):
    """Une question de compréhension pour vérifier l'acquisition."""

    question: str = Field(description="La question posée à l'élève.")
    expected_answer: str = Field(
        description="La réponse attendue, en une à deux phrases."
    )


class Explanation(BaseModel):
    """Réponse pédagogique complète sur un concept Python."""

    concept: str = Field(description="Le nom du concept expliqué.")
    intuition: str = Field(
        description=(
            "Explication intuitive en exactement 3 phrases. "
            "Évite le jargon, vise la clarté pour l'élève."
        )
    )
    code_example: str = Field(
        description=(
            "Un exemple de code Python minimal et commenté qui illustre "
            "le concept. Préfère des noms de variables parlants."
        )
    )
    pitfall: str = Field(
        description="Un piège classique que rencontrent les élèves sur ce concept."
    )
    comprehension_questions: list[ComprehensionQuestion] = Field(
        description="Exactement 3 questions de compréhension, de difficulté croissante."
    )


# Utilisation dans une chaîne LCEL
def explain_chain_for_level(level: str):
    model = get_model().with_structured_output(Explanation)
    return EXPLAIN_PROMPT.partial(level=level) | model
\end{pythoncode}

\subsubsection{Explications pas à pas}

\subsubsection*{Étape 1 — Pourquoi structurer la sortie}

Sans schéma, un LLM répond en texte libre. Le code applicatif doit alors
\emph{parser} cette réponse~: regex, \texttt{json.loads}, conventions de
markdown\dots\ Chaque appel devient un pari sur le fait que le modèle
respectera bien le format demandé. En production, ce pari échoue
régulièrement~: virgule manquante, guillemets inversés, balise oubliée.

\begin{quote}\itshape
\textbf{Faire écrire le format au modèle est moins fiable que le lui
imposer.} La sortie structurée renverse la responsabilité~: ce n'est plus
le modèle qui « doit essayer de respecter » un format ; c'est l'API qui
\textbf{refuse} toute sortie qui ne le respecte pas. Le code applicatif
récupère un objet typé, point.
\end{quote}

\subsubsection*{Étape 2 — Le schéma comme \emph{contrat} de réponse}

Un schéma Pydantic n'est pas un simple type~: c'est une description
\textbf{lisible par le LLM}. Quand on appelle
\texttt{with\_structured\_output(Explanation)}, LangChain transforme la
classe en une \emph{spécification de tool} qui contient~:

\begin{itemize}[leftmargin=*]
  \item Le nom du schéma (\texttt{Explanation}).
  \item Le \texttt{docstring} de la classe (« Réponse pédagogique complète
        sur un concept Python »).
  \item Pour chaque champ~: son nom, son type, et surtout sa
        \texttt{description=}.
\end{itemize}

\begin{quote}\itshape
La \texttt{description=} de chaque \texttt{Field} \textbf{est une partie du
prompt}. Le modèle la lit pour décider quoi mettre dans le champ. Soigne-la
au moins autant que ton message système.
\end{quote}

\subsubsection*{Étape 3 — Les types qui parlent au modèle}

\begin{center}
\begin{tabular}{ll}
\texttt{str}                     & champ texte libre \\
\texttt{int}, \texttt{float}     & nombre, validation automatique \\
\texttt{bool}                    & vrai / faux \\
\texttt{Literal["a", "b", "c"]}  & ensemble \textbf{fermé} de valeurs (= énumération) \\
\texttt{list[X]}                 & liste typée (peut être un sous-schéma) \\
\texttt{X | None}                & champ optionnel \\
\texttt{BaseModel} imbriqué      & sous-objet typé (cf. \texttt{ComprehensionQuestion}) \\
\end{tabular}
\end{center}

Tous ces types sont traduits en JSON Schema. Le LLM reçoit ce schéma et
sait, par exemple, qu'un champ \texttt{Literal} n'a que 3 valeurs
possibles — il ne pourra \emph{physiquement pas} en renvoyer une quatrième.

\subsubsection*{Étape 4 — L'imbrication de schémas}

\begin{center}
\begin{tikzpicture}[node distance=4mm and 10mm]
  \node[schema] (parent) {\texttt{Explanation}};
  \node[schema, below right=8mm and 8mm of parent, xshift=10mm]
      (child)  {\texttt{ComprehensionQuestion}};
  \node[data, right=of parent] (fields)
      {\texttt{concept: str}\\
       \texttt{intuition: str}\\
       \texttt{code\_example: str}\\
       \texttt{pitfall: str}\\
       \texttt{comprehension\_questions:}\\
       \texttt{\ \ list[ComprehensionQuestion]}};
  \node[data, right=of child] (subfields)
      {\texttt{question: str}\\
       \texttt{expected\_answer: str}};

  \draw[flow] (parent) -- (fields);
  \draw[flow] (child)  -- (subfields);
  \draw[flowdata] (parent) -- (child)
       node[label, midway, right, xshift=2mm]{liste imbriquée};
\end{tikzpicture}
\end{center}

\texttt{Explanation} contient une \texttt{list[ComprehensionQuestion]}.
Chaque question dans la liste est elle-même un objet typé, qu'on peut
itérer en Python~:

\begin{pythoncode}
for q in result.comprehension_questions:
    print(q.question)            # str
    print(q.expected_answer)     # str
\end{pythoncode}

\begin{quote}\itshape
\textbf{Les sous-schémas sont la voie idiomatique pour modéliser des listes
d'objets} (questions, étapes, sources, citations\dots). Évite de stocker du
JSON sérialisé dans un champ \texttt{str} — tu perds la validation et
l'auto-complétion.
\end{quote}

\subsubsection*{Étape 5 — Ce qui se passe sous le capot (tool-calling)}

\begin{center}
\begin{tikzpicture}[node distance=6mm and 14mm]
  \node[block]                (a) {Schéma \texttt{Explanation}};
  \node[block, below=of a]    (b) {JSON Schema généré};
  \node[block, below=of b]    (c) {Tool envoyé à l'API Anthropic};
  \node[llm,   right=of c]    (d) {Claude\\(tool-calling)};
  \node[block, above=of d]    (e) {Tool call JSON\\(arguments du tool)};
  \node[block, above=of e]    (f) {Désérialisation\\$\to$ objet \texttt{Explanation}};

  \draw[flow] (a) -- (b) node[label, midway, right]{Pydantic};
  \draw[flow] (b) -- (c) node[label, midway, right]{LangChain};
  \draw[flow] (c) -- (d);
  \draw[flow] (d) -- (e) node[label, midway, right]{Claude « appelle » le tool};
  \draw[flow] (e) -- (f);
\end{tikzpicture}
\end{center}

\texttt{with\_structured\_output} \textbf{n'invente pas} un parseur de
réponse texte. Il déclare un \emph{tool} virtuel auprès du modèle, dont les
arguments \textbf{sont} les champs du schéma. Claude est alors quasi-forcé
d'« appeler » ce tool ; les arguments arrivent en JSON valide, validés par
Pydantic, et reconvertis en objet Python.

\begin{quote}\itshape
Si tu utilises un modèle qui \emph{ne supporte pas} le tool-calling natif
(certains modèles open-source), \texttt{with\_structured\_output} retombe
sur un parsing texte plus fragile. Vérifie toujours que le modèle ciblé
supporte les tools — c'est le cas de tous les Claude \texttt{messages}, de
GPT-4, et de Gemini.
\end{quote}

\subsubsection*{Étape 6 — L'objet de sortie côté Python}

\begin{pythoncode}
chain = explain_chain_for_level("intermédiaire")
result: Explanation = chain.invoke({"concept": "les générateurs Python"})

# result est un VRAI objet Python.
# - PyCharm autocomplète result.intuition, result.pitfall, etc.
# - mypy / Pyright valident les accès.
# - On peut le sérialiser avec result.model_dump_json(indent=2).
print(result.intuition)
for q in result.comprehension_questions:
    print(f"- {q.question}")
\end{pythoncode}

\begin{quote}\itshape
Le \texttt{result: Explanation} en annotation de type est important. Sans
elle, l'IDE pense que \texttt{result} est de type \texttt{Any} et tu perds
l'auto-complétion sur les champs. C'est le \emph{cadeau} de la sortie
structurée — sois sûr de l'utiliser.
\end{quote}

\subsection{Vérification dans LangSmith}

\subsubsection*{L'arbre d'exécution attendu}

\begin{center}
\begin{tikzpicture}[every node/.style={font=\footnotesize\ttfamily, anchor=west}]
  \node                              (a) {RunnableSequence};
  \node[below=1mm of a, xshift=4mm]  (b) {ChatPromptTemplate};
  \node[below=1mm of b, xshift=0mm]  (c) {ChatAnthropic  (structured\_output=Explanation)};
  \node[below=1mm of c, xshift=4mm]  (d) {  Input  : List[BaseMessage]};
  \node[below=1mm of d]              (e) {  Output : Explanation(...)};
  \node[below=1mm of e]              (f) {  Tool call : Explanation(\{\dots\})};
\end{tikzpicture}
\end{center}

\subsubsection*{Les trois choses à regarder}

\begin{enumerate}[leftmargin=*]
  \item \textbf{Onglet \emph{Tool calls} du nœud \texttt{ChatAnthropic}}~:
        tu y vois le tool \texttt{Explanation} déclaré, avec son JSON
        Schema complet (les types et les \texttt{description=} de chaque
        champ). C'est la \textbf{preuve} que ton schéma est bien transmis
        au modèle.
  \item \textbf{Onglet \emph{Output}}~: la réponse brute du modèle est un
        tool-call dont les \texttt{args} contiennent les champs du schéma.
        LangChain la transforme automatiquement en objet \texttt{Explanation}.
  \item \textbf{Onglet \emph{Costs}}~: noter que la sortie structurée coûte
        légèrement plus cher en tokens d'entrée (le schéma est envoyé) mais
        souvent moins en tokens de sortie (le modèle ne génère pas de
        préambule du genre « Voici votre réponse en JSON\dots »).
\end{enumerate}

\begin{quote}\itshape
Modifie une \texttt{description=} d'un champ et relance. Tu vois dans
LangSmith que la nouvelle description est immédiatement présente dans le
tool envoyé au LLM, et que le contenu du champ change conformément. C'est
le \emph{prompt engineering} appliqué directement au schéma de sortie.
\end{quote}

\subsection*{Récapitulatif}
{\itshape\sloppy
\textit{Ce qu'il faut retenir :}
\begin{enumerate}[leftmargin=*]
  \item \textbf{Sortie structurée} = on impose au LLM un schéma de
        réponse, on récupère un \emph{objet Python typé}, plus de parsing
        fragile.
  \item \texttt{pydantic.BaseModel} + \texttt{Field(description=\dots)} est
        la voie standard. Les descriptions sont \emph{lues par le modèle},
        elles font partie du prompt.
  \item \texttt{model.with\_structured\_output(Schema)} fonctionne via le
        \textbf{tool-calling} natif des API modernes — donc fiable, sans
        ré-essais ni parseur custom.
  \item Types puissants à connaître~: \texttt{Literal[\dots]} pour les
        énumérations, \texttt{BaseModel} imbriqué pour les sous-objets,
        \texttt{list[Schema]} pour les listes typées.
  \item Annoter \texttt{result: MonSchema} dans le code applicatif pour
        bénéficier de l'auto-complétion et de la vérification statique.
\end{enumerate}
\par}

% #####################################################################
% #####################################################################
\section{LCEL : composition avec \texttt{|} et figement de variables avec \texttt{.partial()}}
% #####################################################################
% #####################################################################

\begin{tcolorbox}[title=Exercice 3. Composition LCEL avec \texttt{|} et \texttt{.partial()}]
On dispose d'un prompt unique \texttt{EXPLAIN\_PROMPT} qui prend deux
variables \texttt{\{level\}} et \texttt{\{concept\}}. On veut produire
\textbf{trois chaînes spécialisées} (débutant, intermédiaire, avancé),
chacune n'attendant plus que \texttt{\{concept\}} en entrée — le niveau
étant déjà figé.

\textbf{Contrainte}~: pas de duplication du prompt, pas de fonction
\texttt{if/elif} qui réécrit la chaîne à chaque appel. La spécialisation
doit se faire \textbf{à la construction}, et chaque chaîne produite doit
rester un \texttt{Runnable} composable, traçable et streamable.

\bigskip
\noindent\textbf{Objectifs pédagogiques}
\begin{itemize}[leftmargin=*]
  \item Comprendre LCEL (\emph{LangChain Expression Language}) comme une
        \textbf{algèbre de composition} de \texttt{Runnable}.
  \item Maîtriser l'opérateur \texttt{|} et le type qu'il produit.
  \item Utiliser \texttt{.partial()} pour figer une variable d'un prompt et
        produire des \textbf{chaînes spécialisées sans dupliquer de code}.
  \item Reconnaître le pattern \emph{fabrique} (\emph{factory}) pour
        construire dynamiquement des chaînes paramétrées.
\end{itemize}
\end{tcolorbox}

\subsection{Solution}

\subsubsection{Architecture}

\subsubsection*{La composition par \texttt{|}}

\begin{center}
\resizebox{\linewidth}{!}{%
\begin{tikzpicture}[node distance=4mm and 8mm]
  \node[data] (in) {\{"concept": "lambda"\}};
  \node[block, right=of in] (prompt) {prompt\\ (Runnable)};
  \node[pipeop, right=of prompt] (op1) {\textbar};
  \node[llm, right=of op1] (model) {model\\ (Runnable)};
  \node[data, right=of model] (out) {Explanation(\dots)};

  \draw[flowdata] (in)     -- (prompt);
  \draw[flow]     (prompt) -- (op1);
  \draw[flow]     (op1)    -- (model);
  \draw[flow]     (model)  -- (out);

  \begin{scope}[on background layer]
    \node[draw=blue!40!black, dashed, rounded corners=3pt, inner sep=3mm,
          fit=(prompt) (op1) (model), label=below:{\footnotesize\itshape un seul Runnable composé}] {};
  \end{scope}
\end{tikzpicture}%
}
\end{center}

\texttt{prompt | model} crée un nouveau \texttt{Runnable} qui chaîne les
deux. La sortie du gauche devient l'entrée du droit. C'est typé, c'est
composable, c'est traçable.

\subsubsection*{Le figement par \texttt{.partial()}}

\begin{center}
\begin{tikzpicture}[node distance=5mm and 14mm]
  \node[block] (orig) {EXPLAIN\_PROMPT\\\scriptsize attend \texttt{\{level, concept\}}};
  \node[block, below=of orig, fill=yellow!15, draw=orange!70!black]
      (partial) {EXPLAIN\_PROMPT.partial(level="débutant")\\\scriptsize attend juste \texttt{\{concept\}}};
  \node[block, right=24mm of orig] (orig2) {EXPLAIN\_PROMPT};
  \node[block, below=of orig2, fill=yellow!15, draw=orange!70!black]
      (partial2) {.partial(level="avancé")\\\scriptsize attend juste \texttt{\{concept\}}};

  \draw[flow] (orig)  -- (partial)  node[label, midway, right]{fige \texttt{level}};
  \draw[flow] (orig2) -- (partial2) node[label, midway, right]{fige \texttt{level}};
\end{tikzpicture}
\end{center}

\texttt{.partial()} ne mute pas le prompt d'origine~: il renvoie une
\textbf{nouvelle} version partiellement renseignée. C'est l'équivalent de
\texttt{functools.partial} pour les prompts.

\subsubsection{Le code complet}

\begin{pythoncode}
from langchain_core.prompts import ChatPromptTemplate
from config import get_model

EXPLAIN_PROMPT = ChatPromptTemplate.from_messages([
    ("system", "Tu es un assistant pédagogique Python. "
               "Niveau actuel : {level}."),
    ("human",  "Explique : {concept}"),
])

def explain_chain_for_level(level: str):
    """Renvoie une chaîne dont le niveau est figé.

    `.partial(level=...)` crée un nouveau prompt avec une variable fixée.
    On peut donc construire des chaînes spécialisées sans dupliquer le template.
    """
    model = get_model().with_structured_output(Explanation)
    return EXPLAIN_PROMPT.partial(level=level) | model


# Trois chaînes spécialisées, un seul template
beginner_chain     = explain_chain_for_level("débutant")
intermediate_chain = explain_chain_for_level("intermédiaire")
advanced_chain     = explain_chain_for_level("avancé")

# Chaque chaîne n'attend plus que {concept}
result = intermediate_chain.invoke({"concept": "les générateurs Python"})
\end{pythoncode}

\subsubsection{Explications pas à pas}

\subsubsection*{Étape 1 — Qu'est-ce qu'un \texttt{Runnable}}

LCEL repose sur une \textbf{interface unique} : tout objet manipulable est
un \texttt{Runnable}. Cette interface garantit~:

\begin{itemize}[leftmargin=*]
  \item \texttt{.invoke(input)} — exécution synchrone, un input.
  \item \texttt{.batch([inputs])} — exécution synchrone, plusieurs inputs.
  \item \texttt{.stream(input)} — itérateur, pour le streaming.
  \item \texttt{.ainvoke}, \texttt{.abatch}, \texttt{.astream} — versions
        asynchrones.
\end{itemize}

\texttt{ChatPromptTemplate}, un modèle, un parseur, un schéma — tous sont
des \texttt{Runnable}. C'est ce qui permet de les composer librement.

\begin{quote}\itshape
\textbf{Un seul protocole pour tout.} C'est l'innovation de LCEL par
rapport aux versions antérieures de LangChain~: au lieu d'avoir des
\texttt{LLMChain}, \texttt{SequentialChain}, \texttt{RouterChain},
\dots, il n'y a qu'\textbf{un} type composable. Le code applicatif n'a plus
à connaître qu'un seul vocabulaire.
\end{quote}

\subsubsection*{Étape 2 — Le \texttt{|} comme sucre syntaxique}

\texttt{a | b} appelle \texttt{a.\_\_or\_\_(b)}, méthode définie sur la
classe de base \texttt{Runnable}. Elle retourne une instance de
\texttt{RunnableSequence(a, b)}. Mentalement, la transformation est~:

\begin{center}
\begin{tikzpicture}[node distance=4mm and 6mm,
                    every node/.style={font=\small\ttfamily}]
  \node (left) {prompt | model};
  \node[right=of left, font=\Large] (eq) {$\equiv$};
  \node[right=of eq] (right) {RunnableSequence(prompt, model)};
\end{tikzpicture}
\end{center}

Et \texttt{a | b | c} compose en chaîne (associatif à gauche) — équivalent
à \texttt{RunnableSequence(a, b, c)}. Quand on \texttt{.invoke(x)} le
résultat, LCEL fait \texttt{c.invoke(b.invoke(a.invoke(x)))}.

\subsubsection*{Étape 3 — Le contrat de types entre maillons}

\begin{center}
\begin{tikzpicture}[node distance=4mm and 14mm]
  \node[block] (a) {prompt};
  \node[block, right=of a] (b) {model};
  \node[data,  above=of a, font=\scriptsize\ttfamily] (in)  {dict\\\{level, concept\}};
  \node[data,  above=of b, font=\scriptsize\ttfamily] (mid) {ChatPromptValue};
  \node[data,  right=of b, font=\scriptsize\ttfamily] (out) {Explanation\\(BaseModel)};

  \draw[flow] (in)  -- (a);
  \draw[flow] (a)   -- (mid);
  \draw[flow] (mid) -- (b);
  \draw[flow] (b)   -- (out);
\end{tikzpicture}
\end{center}

La sortie d'un maillon doit être acceptable par le suivant. Le modèle
accepte indifféremment un \texttt{ChatPromptValue}, une liste de messages
ou un string — c'est ce qui fait que \texttt{prompt | model} marche
naturellement.

\begin{quote}\itshape
Le \texttt{|} ne vérifie pas les types à la composition. Un mauvais
chaînage (par exemple \texttt{model | prompt}) ne plantera qu'au
\texttt{.invoke()}, avec une erreur souvent obscure. Compose toujours dans
l'ordre \emph{entrée $\to$ traitement $\to$ sortie}.
\end{quote}

\subsubsection*{Étape 4 — \texttt{.partial()} : la fonction partielle des prompts}

Un \texttt{ChatPromptTemplate} qui attend deux variables \texttt{\{level\}}
et \texttt{\{concept\}} ressemble à une fonction Python \texttt{f(level,
concept)}. \texttt{partial(level="débutant")} produit une nouvelle
fonction \texttt{g(concept)} où \texttt{level} est figé.

\begin{pythoncode}
# Le prompt d'origine attend les deux variables.
EXPLAIN_PROMPT.input_variables
# >>> ['level', 'concept']

beginner_prompt = EXPLAIN_PROMPT.partial(level="débutant")
beginner_prompt.input_variables
# >>> ['concept']

# On peut désormais invoquer en ne passant que {concept}.
beginner_prompt.invoke({"concept": "lambda"})
\end{pythoncode}

\begin{quote}\itshape
\texttt{.partial()} \textbf{ne mute pas} le prompt original. C'est une
opération \emph{pure} qui renvoie une nouvelle instance. Tu peux donc
appeler \texttt{EXPLAIN\_PROMPT.partial(\dots)} cent fois sans casser le
template de référence.
\end{quote}

\subsubsection*{Étape 5 — Le pattern \emph{fabrique}}

\texttt{explain\_chain\_for\_level(level)} est une \textbf{fabrique de
chaînes}~: on l'appelle avec un paramètre, elle nous renvoie un
\texttt{Runnable} prêt à l'emploi, déjà configuré.

\begin{center}
\begin{tikzpicture}[node distance=5mm and 14mm]
  \node[block]                  (f) {explain\_chain\_for\_level(level)};
  \node[block, right=of f]      (c1) {beginner\_chain\\\scriptsize (level figé)};
  \node[block, below=of c1]     (c2) {intermediate\_chain};
  \node[block, below=of c2]     (c3) {advanced\_chain};

  \draw[flow] (f) -- (c1) node[label, midway, above]{level="débutant"};
  \draw[flow] (f) -- (c2) node[label, midway, above]{level="intermédiaire"};
  \draw[flow] (f) -- (c3) node[label, midway, above]{level="avancé"};
\end{tikzpicture}
\end{center}

Avantages~:

\begin{itemize}[leftmargin=*]
  \item \textbf{Pas de duplication} du prompt (un seul source de vérité).
  \item \textbf{Spécialisation à la construction}, pas à l'appel. Une fois
        construite, la chaîne ne demande plus le niveau à chaque
        \texttt{invoke}.
  \item \textbf{Composabilité conservée}~: les chaînes spécialisées
        peuvent être passées à un \texttt{RunnableBranch}, un
        \texttt{RunnableParallel}, etc.
\end{itemize}

\begin{quote}\itshape
On vient de remplacer trois fonctions différentes ou trois copies du
prompt par \textbf{une fabrique paramétrée}. C'est l'esprit LCEL~:
construire dynamiquement des graphes d'exécution à partir de briques
primitives, plutôt que d'écrire du code impératif.
\end{quote}

\subsection{Vérification dans LangSmith}

\subsubsection*{Que vois-tu dans l'arbre}

Chaque chaîne produite par la fabrique est un \texttt{RunnableSequence}.
Quand tu \texttt{invoke} \texttt{beginner\_chain}, LangSmith trace~:

\begin{center}
\begin{tikzpicture}[every node/.style={font=\footnotesize\ttfamily, anchor=west}]
  \node                              (a) {RunnableSequence (beginner\_chain)};
  \node[below=1mm of a, xshift=4mm]  (b) {ChatPromptTemplate};
  \node[below=1mm of b, xshift=4mm]  (c) {  Input  : \{concept: "lambda"\}};
  \node[below=1mm of c]              (d) {  Output : SystemMessage("\dots Niveau : débutant\dots")};
  \node[below=1mm of d, xshift=-4mm] (e) {ChatAnthropic (structured\_output=Explanation)};
\end{tikzpicture}
\end{center}

Le \texttt{level} \textbf{n'apparaît plus dans l'input} du run — il a été
absorbé par le \texttt{.partial()}. C'est la preuve visuelle que le
figement est bien effectif.

\subsubsection*{Les trois choses à regarder}

\begin{enumerate}[leftmargin=*]
  \item \textbf{Le \emph{type} du run racine}~: \texttt{RunnableSequence}.
        C'est la matérialisation du \texttt{|}.
  \item \textbf{L'input du \texttt{ChatPromptTemplate}}~: seulement la (ou
        les) variable(s) non figée(s). Si tu y vois encore \texttt{level},
        c'est que ton \texttt{.partial()} n'a pas été pris en compte.
  \item \textbf{L'output du \texttt{ChatPromptTemplate}}~: les messages
        complets avec \texttt{level} \emph{déjà remplacé}. La substitution
        a eu lieu à la construction du template, pas dans le modèle.
\end{enumerate}

\begin{quote}\itshape
Le tracing LangSmith permet de \emph{voir} la différence entre
« variable figée par \texttt{.partial()} » et « variable passée à
\texttt{.invoke()} »~: la première est invisible dans l'input du run, la
seconde est visible. C'est utile pour déboguer les chaînes paramétrées.
\end{quote}

\subsection*{Récapitulatif}
{\itshape\sloppy
\textit{Ce qu'il faut retenir :}
\begin{enumerate}[leftmargin=*]
  \item LCEL repose sur \textbf{une seule interface}~: \texttt{Runnable}.
        Tout objet manipulable (prompt, modèle, parseur, branch, parallel,
        \dots) implémente \texttt{.invoke}, \texttt{.batch},
        \texttt{.stream}, et leurs versions async.
  \item \texttt{a | b} crée un \texttt{RunnableSequence(a, b)}. C'est
        purement du sucre syntaxique sur la méthode
        \texttt{Runnable.\_\_or\_\_}.
  \item \texttt{ChatPromptTemplate.partial(\dots)} fige une ou plusieurs
        variables sans muter le template original. La chaîne produite
        n'attend plus que les variables restantes.
  \item \textbf{Pattern fabrique}~: une fonction \texttt{build\_xxx(\dots)}
        qui renvoie un \texttt{Runnable} paramétré est la manière idiomatique
        de spécialiser à la construction plutôt qu'à l'appel.
  \item Dans LangSmith, le \texttt{|} apparaît comme un nœud
        \texttt{RunnableSequence} ; \texttt{.partial()} se voit à
        l'\emph{absence} de la variable figée dans l'input du run.
\end{enumerate}
\par}

% #####################################################################
% #####################################################################
\section{\texttt{RunnableParallel} : exécuter plusieurs chaînes en concurrence}
% #####################################################################
% #####################################################################

\begin{tcolorbox}[title=Exercice 4. \texttt{RunnableParallel}]
On veut générer \textbf{deux versions} d'une même explication~: une pour
les débutants, une pour les intermédiaires. Les deux explications sont
\textbf{indépendantes} l'une de l'autre (aucune n'a besoin du résultat de
l'autre).

\textbf{Contrainte de performance}~: chaque appel au LLM prend
$\sim$1 à 2 secondes. En exécution séquentielle, on attendrait
$\sim$2 à 4 secondes au total. On veut \textbf{ne payer que la latence
d'un seul appel}, soit $\sim$1 à 2 secondes — et un seul \texttt{invoke}
côté code applicatif.

\bigskip
\noindent\textbf{Objectifs pédagogiques}
\begin{itemize}[leftmargin=*]
  \item Comprendre comment LCEL \emph{parallélise} des chaînes indépendantes
        avec \texttt{RunnableParallel}.
  \item Quantifier le gain de latence par rapport à une boucle séquentielle.
  \item Distinguer \texttt{RunnableParallel} de
        \texttt{RunnablePassthrough.assign} (qui paraît proche mais a une
        sémantique différente).
  \item Repérer la structure « en éventail » dans LangSmith.
\end{itemize}
\end{tcolorbox}

\subsection{Solution}

\subsubsection{Architecture}

\subsubsection*{Vue d'ensemble}

\begin{center}
\begin{tikzpicture}
  \node[data] (in) {\{"concept": "générateurs"\}};
  \node[parallel, below=of in] (par) {RunnableParallel\\\texttt{beginner=\dots, intermediate=\dots}};

  \node[llm, below left=12mm and 4mm of par]   (a) {beginner\_chain\\\scriptsize (level="débutant")};
  \node[llm, below right=12mm and 4mm of par]  (b) {intermediate\_chain\\\scriptsize (level="intermédiaire")};

  \node[data, below=22mm of par] (out)
       {\{"beginner": Explanation(\dots),\\\ "intermediate": Explanation(\dots)\}};

  \draw[flowdata] (in)  -- (par);
  \draw[flow]     (par) -- (a) node[label, midway, sloped, above]{en parallèle};
  \draw[flow]     (par) -- (b) node[label, midway, sloped, above]{en parallèle};
  \draw[flow]     (a)   -- (out);
  \draw[flow]     (b)   -- (out);
\end{tikzpicture}
\end{center}

\texttt{RunnableParallel} duplique l'input vers chaque branche, lance toutes
les branches \textbf{en concurrence}, attend que toutes terminent, et
retourne un \textbf{dict} dont les clés sont les noms de branches.

\subsubsection*{Diagramme temporel : séquentiel vs parallèle}

\begin{center}
\begin{tikzpicture}[node distance=3mm and 6mm]
  % Axe séquentiel
  \draw[timearrow] (0,0) -- (10,0) node[right]{\small temps};
  \node[left] at (-0.2, 0.6) {\scriptsize séquentiel};
  \node[bar, anchor=west, minimum width=20mm] (s1) at (0, 0.6) {beginner};
  \node[bar, anchor=west, minimum width=22mm, fill=orange!30] (s2) at (2, 0.6) {intermediate};

  \draw[decorate, decoration={brace, amplitude=2mm}]
       (0, 1.2) -- (4.2, 1.2) node[midway, above, yshift=2mm]{\scriptsize $\sim$4 s};

  % Axe parallèle
  \draw[timearrow] (0,-2) -- (10,-2);
  \node[left] at (-0.2, -1.4) {\scriptsize parallèle};
  \node[bar, anchor=west, minimum width=20mm] (p1) at (0, -1.4) {beginner};
  \node[bar, anchor=west, minimum width=22mm, fill=orange!30] (p2) at (0, -1.95) {intermediate};

  \draw[decorate, decoration={brace, amplitude=2mm}]
       (0, -1.0) -- (2.2, -1.0) node[midway, above, yshift=2mm]{\scriptsize $\sim$2 s};
\end{tikzpicture}
\end{center}

Le gain est d'autant plus marqué qu'on a de branches~: 3 chaînes en
parallèle restent à la latence d'\emph{une} seule chaîne (modulo le
plafond de \emph{rate limiting} du fournisseur).

\subsubsection{Le code complet}

\begin{pythoncode}
from langchain_core.runnables import RunnableParallel

def explain_chain_for_level(level: str):
    model = get_model().with_structured_output(Explanation)
    return EXPLAIN_PROMPT.partial(level=level) | model


def build_multi_level_chain():
    """Génère les versions débutant ET intermédiaire en un seul invoke."""
    return RunnableParallel(
        beginner=explain_chain_for_level("débutant"),
        intermediate=explain_chain_for_level("intermédiaire"),
    )


# Utilisation
multi = build_multi_level_chain()
results = multi.invoke({"concept": "les générateurs Python"})

# results est un dict.
# Les deux appels LLM ont été lancés en concurrence ; on n'a payé
# que la latence du plus lent.
print(results["beginner"].intuition)
print(results["intermediate"].intuition)
\end{pythoncode}

\subsubsection{Explications pas à pas}

\subsubsection*{Étape 1 — Le contrat de \texttt{RunnableParallel}}

\texttt{RunnableParallel} prend un \textbf{ensemble de} \texttt{Runnable}
\textbf{nommés}, généralement passés en kwargs~:

\begin{pythoncode}
RunnableParallel(
    beginner=chain_a,
    intermediate=chain_b,
    advanced=chain_c,
)
\end{pythoncode}

À l'invocation~:

\begin{enumerate}[leftmargin=*]
  \item L'input \texttt{x} est \textbf{dupliqué} vers chaque branche.
  \item Toutes les branches s'exécutent en parallèle (via le pool de
        threads pour les chaînes synchrones, via l'event loop asyncio
        pour les chaînes async).
  \item LCEL attend que \textbf{toutes} terminent.
  \item La sortie est un dict \texttt{\{nom\_branche: résultat\}}.
\end{enumerate}

\begin{quote}\itshape
\textbf{Un seul input, plusieurs sorties.} C'est le pattern « éventail »~:
on prend une question, on en demande plusieurs angles, et chaque branche
travaille sur sa version.
\end{quote}

\subsubsection*{Étape 2 — Le gain de latence n'est pas le gain de \emph{débit}}

Un appel LLM passe la quasi-totalité du temps à \textbf{attendre} la
réponse du serveur. Sur le client, le thread est endormi. Lancer deux
appels en parallèle ne consomme donc pas deux fois plus de CPU local~: ça
consomme deux \emph{tickets} chez le fournisseur en même temps.

\begin{center}
\begin{tabular}{lll}
\toprule
\textbf{Mesure} & \textbf{Séquentiel} & \textbf{Parallèle (2 branches)} \\
\midrule
Latence totale & $T_1 + T_2$        & $\max(T_1, T_2)$ \\
Tokens consommés & $N_1 + N_2$      & $N_1 + N_2$ \\
Coût \$         & idem               & idem \\
CPU client      & 1 thread occupé   & 2 threads endormis sur I/O \\
\bottomrule
\end{tabular}
\end{center}

\begin{quote}\itshape
La parallélisation \textbf{ne réduit pas la facture}~: tu paies toujours
chaque token. Elle réduit uniquement la \emph{latence perçue} par
l'utilisateur. Pour un usage UI où le temps de réponse compte, c'est
décisif. Pour un job batch nocturne, le gain est nul (et tu risques
juste de saturer plus vite ton \emph{rate limit}).
\end{quote}

\subsubsection*{Étape 3 — Le sucre syntaxique dict}

\texttt{RunnableParallel} a deux formes équivalentes~:

\begin{pythoncode}
# Forme explicite
chain = RunnableParallel(beginner=a, intermediate=b)

# Forme dict (utile pour des noms non-identifieurs, ex. avec espaces)
chain = RunnableParallel({"beginner": a, "intermediate": b})
\end{pythoncode}

Et LCEL accepte même un \textbf{dict nu} dans une composition par
\texttt{|}~: il sera implicitement converti en \texttt{RunnableParallel}.

\begin{pythoncode}
chain = some_prompt | {"answer": chain_a, "summary": chain_b}
# strictement équivalent à
chain = some_prompt | RunnableParallel(answer=chain_a, summary=chain_b)
\end{pythoncode}

\begin{quote}\itshape
Le \textbf{dict implicite} dans une chaîne LCEL signifie toujours
\texttt{RunnableParallel}. Quand tu vois \texttt{\{ "x": f, "y": g \}} dans
un \texttt{|}, lis « éventail vers f et g, sortie nommée ».
\end{quote}

\subsubsection*{Étape 4 — \texttt{RunnableParallel} vs \texttt{RunnablePassthrough.assign}}

Ces deux primitives se ressemblent, mais l'\emph{intention} diffère~:

\begin{center}
\begin{tabular}{p{6.5cm} p{6.5cm}}
\toprule
\textbf{\texttt{RunnableParallel}} & \textbf{\texttt{RunnablePassthrough.assign}} \\
\midrule
Renvoie un \textbf{nouveau dict} contenant uniquement les clés des branches.
& \textbf{Enrichit} l'input existant~: garde toutes les clés d'entrée \emph{et} ajoute les nouvelles.
\\[1mm]
Sortie : \texttt{\{nom1: \dots, nom2: \dots\}} (les inputs disparaissent).
& Sortie : \texttt{\{**input, nom1: \dots, nom2: \dots\}} (inputs préservés).
\\[1mm]
Usage typique~: fan-out final (l'utilisateur veut \emph{le résultat}).
& Usage typique~: ajout d'un champ calculé à mi-pipeline.
\\
\bottomrule
\end{tabular}
\end{center}

\begin{quote}\itshape
Si tu utilises \texttt{RunnableParallel} à la place de
\texttt{.assign(\dots)} dans un pipeline, l'\texttt{input} d'origine
\textbf{disparaît} dans la sortie. Les maillons suivants qui s'attendaient
à lire la clé \texttt{concept} ne la trouveront plus. C'est l'un des bugs
LCEL les plus fréquents.
\end{quote}

\subsubsection*{Étape 5 — Limites et bonnes pratiques}

\begin{itemize}[leftmargin=*]
  \item \textbf{Rate limits}~: chaque branche compte. 5 branches en
        parallèle, c'est 5 appels simultanés chez le fournisseur — vérifie
        que tu ne dépasses pas le RPM (\emph{requests per minute}) de ton
        compte.
  \item \textbf{Erreur partielle}~: si une branche échoue, par défaut tout
        échoue. Pour tolérer les erreurs partielles, enveloppe chaque
        branche dans \texttt{.with\_fallbacks([\dots])}.
  \item \textbf{Branches indépendantes uniquement}~: si la branche B a
        besoin du résultat de la branche A, ce n'est pas du parallèle —
        c'est une séquence, à composer avec \texttt{|} et \texttt{.assign}.
\end{itemize}

\subsubsection{Exécution et observation}

\subsubsection*{Le code de l'exécuteur}

\begin{pythoncode}
import time

multi = build_multi_level_chain()

t0 = time.perf_counter()
results = multi.invoke({"concept": "les générateurs Python"})
elapsed = time.perf_counter() - t0

print(f"Durée RunnableParallel : {elapsed:.2f} s")

# Pour comparaison, version séquentielle
t0 = time.perf_counter()
r1 = explain_chain_for_level("débutant").invoke({"concept": "les générateurs Python"})
r2 = explain_chain_for_level("intermédiaire").invoke({"concept": "les générateurs Python"})
elapsed_seq = time.perf_counter() - t0

print(f"Durée séquentielle : {elapsed_seq:.2f} s")
print(f"Speedup : x{elapsed_seq / elapsed:.2f}")
\end{pythoncode}

Typiquement, on observe un speedup proche de $\times 1.8$ à $\times 2.0$
pour deux branches (le facteur 2 théorique étant légèrement entamé par
l'overhead de \emph{threading} et les variations de latence côté
fournisseur).

\subsection{Vérification dans LangSmith}

\subsubsection*{L'arbre d'exécution attendu}

\begin{center}
\begin{tikzpicture}[every node/.style={font=\footnotesize\ttfamily, anchor=west}]
  \node                              (a) {RunnableParallel<beginner, intermediate>};
  \node[below=1mm of a, xshift=4mm]  (b) {[beginner] RunnableSequence};
  \node[below=1mm of b, xshift=4mm]  (c) {  ChatPromptTemplate};
  \node[below=1mm of c]              (d) {  ChatAnthropic (structured\_output) ← LLM \#1};
  \node[below=1mm of d, xshift=-4mm] (e) {[intermediate] RunnableSequence};
  \node[below=1mm of e, xshift=4mm]  (f) {  ChatPromptTemplate};
  \node[below=1mm of f]              (g) {  ChatAnthropic (structured\_output) ← LLM \#2};
\end{tikzpicture}
\end{center}

Les deux sous-séquences sont des \textbf{frères} sous le même
\texttt{RunnableParallel}. LangSmith affiche un \emph{timeline} où l'on
voit clairement les deux barres se chevaucher temporellement.

\subsubsection*{Les quatre choses à regarder}

\begin{enumerate}[leftmargin=*]
  \item \textbf{Timeline en haut du run racine}~: deux barres parallèles.
        Si les barres sont en \emph{série} (l'une après l'autre), c'est
        que \texttt{RunnableParallel} n'a pas été utilisé correctement
        (par exemple, tu as fait deux \texttt{invoke} séparés dans une
        boucle).
  \item \textbf{Input des deux branches}~: identique. C'est la duplication
        de l'input qu'on a évoquée.
  \item \textbf{Output du run racine}~: un dict avec les clés
        \texttt{beginner} et \texttt{intermediate}. Aucune autre clé.
  \item \textbf{Onglet \emph{Costs}}~: les tokens et le \$ s'additionnent
        entre branches — la parallélisation ne change rien à la facture.
\end{enumerate}

\begin{quote}\itshape
Lancer deux appels LLM en parallèle dans LangChain, c'est littéralement
\textbf{ajouter un mot}~: \texttt{RunnableParallel}. Pas de threads à
gérer, pas d'\texttt{asyncio.gather}, pas de pool. Le runtime LCEL fait
tout, et le tracing LangSmith te montre la concurrence à l'œil nu.
\end{quote}

\subsection*{Récapitulatif}
{\itshape\sloppy
\textit{Ce qu'il faut retenir :}
\begin{enumerate}[leftmargin=*]
  \item \texttt{RunnableParallel(nom1=chain1, nom2=chain2, \dots)} duplique
        l'input vers chaque branche, exécute toutes les branches en
        \textbf{concurrence}, et renvoie un \textbf{dict} de résultats.
  \item Latence~: $\max(T_i)$ au lieu de $\sum T_i$. Coût en tokens
        \textbf{inchangé}. C'est un gain de latence, pas un gain de débit.
  \item \texttt{prompt | \{"a": chain\_a, "b": chain\_b\}} est du sucre
        syntaxique~: le dict est implicitement converti en
        \texttt{RunnableParallel}.
  \item \textbf{Ne pas confondre} avec \texttt{RunnablePassthrough.assign}~:
        \texttt{Parallel} remplace l'input, \texttt{.assign} l'enrichit.
  \item Dans LangSmith, la concurrence se voit immédiatement~: deux barres
        de timeline qui se chevauchent sous le même nœud parent.
\end{enumerate}
\par}

% #####################################################################
% #####################################################################
\section{Streaming token par token avec \texttt{.stream()}}
% #####################################################################
% #####################################################################

\begin{tcolorbox}[title=Exercice 5. Streaming token par token]
L'utilisateur pose une question au tuteur Python. Plutôt que d'attendre
\emph{toute} la réponse (souvent 2 à 5 secondes), on veut afficher la
réponse \textbf{au fur et à mesure} qu'elle est générée, comme dans
ChatGPT~: le premier mot apparaît au bout de $\sim$300 ms, puis les
suivants s'écrivent en temps réel.

\textbf{Contrainte}~: rester un \texttt{Runnable} LCEL classique, et
afficher token par token dans le terminal, sans buffering perturbateur.

\bigskip
\noindent\textbf{Objectifs pédagogiques}
\begin{itemize}[leftmargin=*]
  \item Comprendre la différence entre \texttt{.invoke()} (réponse unique)
        et \texttt{.stream()} (itérateur de morceaux).
  \item Savoir construire une chaîne \emph{streamable} et l'afficher au fil
        de l'eau dans un terminal ou une UI.
  \item Comprendre pourquoi \texttt{with\_structured\_output} \emph{ne
        streame pas bien} et quand basculer sur \texttt{StrOutputParser}.
  \item Connaître les pièges classiques du streaming (buffering, flush,
        retour chariot).
\end{itemize}
\end{tcolorbox}

\subsection{Solution}

\subsubsection{Architecture}

\subsubsection*{Vue d'ensemble}

\begin{center}
\resizebox{\linewidth}{!}{%
\begin{tikzpicture}
  \node[data]               (in)     {\{"concept": "le GIL"\}};
  \node[block, below=of in] (prompt) {ChatPromptTemplate};
  \node[llm,   below=of prompt] (model) {Modèle de chat};
  \node[block, below=of model] (parser) {StrOutputParser};

  \node[token, right=22mm of model, yshift=4mm]   (t1) {"Le"};
  \node[token, right=of t1]   (t2) {" GIL"};
  \node[token, right=of t2]   (t3) {" est"};
  \node[token, right=of t3]   (t4) {" un"};
  \node[token, right=of t4]   (t5) {"\dots"};

  \draw[flowdata] (in)     -- (prompt);
  \draw[flow]     (prompt) -- (model);
  \draw[flow]     (model)  -- (parser);
  \draw[flow]     (model.east) -- ++(0.5,0)  node[label, above]{streaming};
  \draw[->, thick, draw=teal!60!black] (model.east) ..controls +(1,0.5).. (t1.west);

  \node[font=\footnotesize\itshape, below=12mm of parser]
       (out) {pour chunk in chain.stream(\dots): print(chunk, end="", flush=True)};
\end{tikzpicture}%
}
\end{center}

L'API du modèle ouvre une connexion \emph{server-sent events} avec le
fournisseur. Chaque token reçu est immédiatement émis par la chaîne. Tant
qu'on consomme l'itérateur, les tokens arrivent.

\subsubsection{Le code complet}

\begin{pythoncode}
from langchain_core.prompts import ChatPromptTemplate
from langchain_core.output_parsers import StrOutputParser
from config import get_model

STREAM_PROMPT = ChatPromptTemplate.from_messages([
    ("system", "Tu es un assistant pédagogique Python concis et précis."),
    ("human",  "Explique en 2 paragraphes : {concept}"),
])


def build_streaming_chain():
    """Chaîne texte simple, idéale pour du streaming dans un terminal/UI."""
    return STREAM_PROMPT | get_model() | StrOutputParser()


# Utilisation
streamer = build_streaming_chain()
for chunk in streamer.stream({"concept": "le GIL en Python"}):
    print(chunk, end="", flush=True)
print()  # saut de ligne final
\end{pythoncode}

\subsubsection{Explications pas à pas}

\subsubsection*{Étape 1 — \texttt{.invoke()} vs \texttt{.stream()}}

\begin{center}
\begin{tabular}{p{6.5cm} p{6.5cm}}
\toprule
\textbf{\texttt{.invoke(input)}} & \textbf{\texttt{.stream(input)}} \\
\midrule
Renvoie une valeur unique~: la réponse complète. & Renvoie un \emph{itérateur} de morceaux (\texttt{chunks}). \\
Le code attend que tout soit prêt avant de continuer. & Le code itère et reçoit les morceaux dès qu'ils arrivent. \\
Latence perçue = latence totale. & Latence perçue = latence du \emph{premier} token (~300 ms). \\
Bien pour batch, scripts, données structurées. & Bien pour UI temps réel, terminal interactif. \\
\bottomrule
\end{tabular}
\end{center}

\begin{quote}\itshape
Le \textbf{time-to-first-token (TTFT)} est ce qui rend une interface
« vivante ». Même si la durée totale est identique, voir la réponse
\textbf{commencer} à apparaître au bout de 300 ms change radicalement la
sensation d'attente.
\end{quote}

\subsubsection*{Étape 2 — Pourquoi \texttt{StrOutputParser} ?}

Par défaut, un modèle de chat renvoie des \texttt{AIMessageChunk} (des
fragments d'\texttt{AIMessage}). En streaming, on ne veut pas avoir à
faire \texttt{chunk.content} dans la boucle d'affichage — on veut
directement le \emph{texte}. C'est exactement le rôle de
\texttt{StrOutputParser}~:

\begin{center}
\begin{tabular}{ll}
Entrée & Sortie \\
\midrule
\texttt{AIMessageChunk(content="Le")}  & \texttt{"Le"} \\
\texttt{AIMessageChunk(content=" GIL")} & \texttt{" GIL"} \\
\dots & \dots \\
\end{tabular}
\end{center}

\begin{pythoncode}
# Sans StrOutputParser
for chunk in (STREAM_PROMPT | get_model()).stream({"concept": "X"}):
    print(chunk.content, end="", flush=True)  # accès manuel à .content

# Avec StrOutputParser
for chunk in (STREAM_PROMPT | get_model() | StrOutputParser()).stream({"concept": "X"}):
    print(chunk, end="", flush=True)  # le chunk EST une str
\end{pythoncode}

\subsubsection*{Étape 3 — Pourquoi \texttt{with\_structured\_output} streame mal}

\begin{center}
\resizebox{\linewidth}{!}{%
\begin{tikzpicture}[node distance=4mm and 8mm]
  \node[block]                          (a) {Schéma Pydantic};
  \node[llm,   right=of a]              (b) {LLM (tool-calling)};
  \node[block, right=of b]              (c) {Désérialisation\\Pydantic};
  \node[data,  right=of c, fill=red!10] (d) {Objet \emph{valide} ?};

  \draw[flow] (a) -- (b);
  \draw[flow] (b) -- (c) node[label, midway, above]{JSON partiel};
  \draw[flow] (c) -- (d) node[label, midway, above]{?};
\end{tikzpicture}%
}
\end{center}

Un schéma Pydantic est \textbf{tout ou rien}~: tant que tous les champs
requis ne sont pas remplis, l'objet ne peut pas être validé. Le modèle
peut bien streamer le JSON token par token, on ne peut pas
\emph{désérialiser} en cours de route — la moitié des champs manquent.

\begin{quote}\itshape
Appeler \texttt{.stream()} sur une chaîne \texttt{with\_structured\_output}
ne plantera \emph{pas}, mais souvent la chaîne accumule tout en interne et
n'émet qu'\textbf{un seul gros chunk} à la fin. Le bénéfice du streaming
disparaît. Pour un streaming réel, soit on accepte du texte
(\texttt{StrOutputParser}), soit on utilise des techniques avancées
(parsing JSON incrémental, mais c'est une autre histoire).
\end{quote}

\subsubsection*{Étape 4 — La boucle d'affichage côté Python}

\begin{pythoncode}
for chunk in streamer.stream({"concept": "le GIL en Python"}):
    print(chunk, end="", flush=True)
print()
\end{pythoncode}

Trois détails subtils mais importants~:

\begin{itemize}[leftmargin=*]
  \item \texttt{end=""}~: empêche \texttt{print} d'ajouter un \texttt{\textbackslash n} après
        chaque chunk. Sans ça, chaque token est affiché sur sa propre
        ligne — illisible.
  \item \texttt{flush=True}~: force l'écriture sur le terminal
        \textbf{immédiatement}. Python bufferise stdout par défaut (par
        bloc de 4 ko ou par ligne) ; sans flush, on voit la réponse arriver
        par paquets, pas par tokens.
  \item Le \texttt{print()} final est juste là pour le saut de ligne après
        toute la réponse.
\end{itemize}

\begin{quote}\itshape
\textbf{\texttt{flush=True} est le secret du streaming visible.} Si tu
oublies, tu as l'impression que ton streaming ne marche pas — alors qu'il
marche, mais c'est le terminal qui retient les caractères. C'est la
question \#1 posée en debug.
\end{quote}

\subsubsection*{Étape 5 — Cas async : \texttt{.astream()}}

Pour une UI web ou un endpoint FastAPI, on veut la version asynchrone~:

\begin{pythoncode}
async def stream_to_websocket(websocket, concept: str):
    async for chunk in streamer.astream({"concept": concept}):
        await websocket.send_text(chunk)
\end{pythoncode}

La sémantique est identique à \texttt{.stream()}, mais on utilise
\texttt{async for} et chaque \texttt{await} libère l'event loop pour les
autres connexions. C'est le pattern à connaître pour servir N clients
simultanément.

\subsection{Vérification dans LangSmith}

\subsubsection*{L'arbre d'exécution attendu}

\begin{center}
\begin{tikzpicture}[every node/.style={font=\footnotesize\ttfamily, anchor=west}]
  \node                              (a) {RunnableSequence};
  \node[below=1mm of a, xshift=4mm]  (b) {ChatPromptTemplate};
  \node[below=1mm of b, xshift=0mm]  (c) {ChatAnthropic  (streaming=True)};
  \node[below=1mm of c, xshift=4mm]  (d) {  Tokens : "Le", " GIL", " est", " un", ...};
  \node[below=1mm of d, xshift=-4mm] (e) {StrOutputParser};
\end{tikzpicture}
\end{center}

LangSmith reconstitue la réponse complète à la fin et l'affiche dans
l'onglet \emph{Output}, mais conserve la durée \emph{end-to-end} et le
\emph{TTFT} dans les métriques.

\subsubsection*{Les trois choses à regarder}

\begin{enumerate}[leftmargin=*]
  \item \textbf{Métrique \emph{Time to first token}}~: temps entre l'appel
        et le premier token reçu. C'est la métrique qui détermine la
        \emph{sensation de vivacité} de ton UI. Souvent affichée sous
        \texttt{ttft\_ms} dans les détails du run.
  \item \textbf{Onglet \emph{Output} du nœud \texttt{ChatAnthropic}}~:
        l'\texttt{AIMessage} reconstitué (LangSmith concatène les chunks
        pour l'affichage). Tu ne vois \emph{pas} chaque token séparément
        dans l'UI, mais la trace HTTP sous-jacente est bien du streaming.
  \item \textbf{Type du modèle dans la trace}~:
        \texttt{streaming=True} ou un flag équivalent. Si tu vois
        \texttt{streaming=False}, ton appel est en mode bloquant.
\end{enumerate}

\begin{quote}\itshape
Compare deux runs côte à côte~: l'un avec \texttt{.invoke}, l'autre avec
\texttt{.stream}, sur le même prompt. Le total de tokens et le coût \$
sont identiques. Mais le run streamé a un TTFT $\sim$10$\times$ plus
court. C'est, mesurée, la valeur ajoutée du streaming.
\end{quote}

\subsection*{Récapitulatif}
{\itshape\sloppy
\textit{Ce qu'il faut retenir :}
\begin{enumerate}[leftmargin=*]
  \item \texttt{chain.stream(input)} renvoie un \textbf{itérateur} de
        morceaux ; \texttt{chain.invoke(input)} renvoie la réponse
        complète. Même chaîne, deux modes de consommation.
  \item Pour streamer du texte, terminer la chaîne par
        \texttt{StrOutputParser} pour que chaque morceau soit déjà une
        \texttt{str} (et pas un \texttt{AIMessageChunk}).
  \item \texttt{with\_structured\_output} ne streame pas \emph{vraiment}~:
        un schéma est tout-ou-rien, on ne peut pas désérialiser un objet
        partiel. Pour du streaming réel, accepter du texte.
  \item Dans la boucle d'affichage~: \texttt{print(chunk, end="",
        flush=True)}. Sans \texttt{flush=True}, les buffers du terminal
        masquent l'effet streaming.
  \item Pour le web ou les apps async, utiliser \texttt{.astream()} avec
        \texttt{async for}. Sémantique identique, mais l'event loop reste
        libre pour les autres connexions.
\end{enumerate}
\par}

\end{document}